\section{Past Actions of the House}

\subsection{The U.S. Commercial Space Launch Competitiveness Act (H.R. 2262, 2015)}

Also known as the SPACE Act of 2015 (Spurring Private Aerospace Competitiveness and Entrepreneurship Act), this is one of the most significant pieces of legislation in this area.\footnote{U.S. Congress. (2015). \emph{H.R.2262 — U.S. Commercial Space Launch Competitiveness Act, 114th Congress (2015–2016)}. Congress.gov. \url{https://www.congress.gov/bill/114th-congress/house-bill/2262}} The bill passed the House and was signed into law by President Obama. It was described as “one of the most significant modernisations of commercial space policy and regulatory legislation since the original Commercial Space Launch Act was enacted in 1984”.\footnote{U.S. House Committee on Science, Space, and Technology. (2015). \emph{H.R. 2262, U.S. Commercial Space Launch Competitiveness Act: What they're saying} [One-pager]. \url{https://science.house.gov/_cache/files/e/0/e05fce77-6f88-4b19-8fba-dcf1c81c3d78/20ECF71186C40C11147C915FC5DB563DDA4BAF396C9E9C274CC7A2CDC7414FD5.h.r.-2262---what-they-re-saying.pdf}}

Crucially, it gave corporations, for the first time, a right to the resources they extract from celestial bodies, providing “legal certainty regarding property rights in asteroid resources”.\footnote{Virginia Law Review. (n.d.). \emph{Mining for meaning: An examination of the legality of property rights in space resources}. \url{https://virginialawreview.org/articles/mining-meaning-examination-legality-property-rights-space-resources/}} The Act also extended liability and insurance protections for spaceflight participants and addressed space traffic management of federal and private assets in outer space.

\subsection{The American Space Commerce Free Enterprise Act (H.R. 2809, 2017)}

Introduced in the 115th Congress by House Science, Space, and Technology Committee Chairman Lamar Smith, this bill was unanimously approved by the Science Committee on 8 June 2017 and subsequently passed the full House unanimously on 24 April 2018.\footnote{U.S. House Committee on Science, Space, and Technology. (2018, April 24). \emph{House approves American Space Commerce Free Enterprise Act} [Press release]. \url{https://science.house.gov/2018/4/house-approves-american-space-commerce-free-enterprise-act}}

The bill notably declared that “space is not a global commons,” a significant departure from the principles underlying ratified international treaties, and directed the President to “protect the interests of U.S. entities in outer space, including commercial activity, the exploitation of space resources, and ownership rights over space objects and obtained space resources”.\footnote{Tronchetti, F. (2017). The American Space Commerce Free Enterprise Act of 2017: The latest step in the US commercial space legislation. \emph{Space Policy}. \url{https://www.sciencedirect.com/science/article/abs/pii/S0265964617301261}}

It expanded the authorities of the Office of Space Commerce within the Department of Commerce to supervise commercial space activity, established a certification process for the private sector to operate objects in outer space, and created a “Private Space Activity Advisory Committee” to advise on U.S. private-sector activities in outer space.\footnote{U.S. Congress. (2018). \emph{Text — H.R.2809 — American Space Commerce Free Enterprise Act, 115th Congress (2017–2018)}. Congress.gov. \url{https://www.congress.gov/bill/115th-congress/house-bill/2809/text}} The bill was framed as declaring that “America is fully ‘open for business’ in space.”\footnote{Footnote: U.S. House Committee on Science, Space, and Technology. (2018, April 24). \emph{House approves American Space Commerce Free Enterprise Act} [Press release]. \url{https://science.house.gov/2018/4/house-approves-american-space-commerce-free-enterprise-act}}

\subsection{The Space Commercialisation Promotion Act (1996)}

Earlier efforts also emerged in the House. The Space Commercialisation Promotion Act of 1996 sought to promote the economic development of outer space through government-created technologies transitioning to the private sector, emphasising the importance of “unleashing the spirit, energy, and work ethic of U.S. entrepreneurialism” to open up the potential of space. \footnote{U.S. House Committee on Science. (1996). \emph{H. Rept. 104-801 — Space Commercialization Promotion Act of 1996}. Congress.gov. \url{https://www.congress.gov/committee-report/104th-congress/house-report/801}}

\subsection{The Commercial Space Act of 2023}

On 2 November 2023, Representatives Frank Lucas (R-OK), as House Science Committee chairman, and Brian Babin (R-TX), as Space and Aeronautics Subcommittee chairman, introduced the Commercial Space Act of 2023, which proposed establishing a certification process for private spacecraft and space objects not currently licensed by the U.S. government.\footnote{Center for Strategic and International Studies, Aerospace Security Project. (n.d.). \emph{Mission authorization: Decoding the space policy dilemma}. \url{https://aerospace.csis.org/mission-authorization-decoding-the-space-policy-dilemma/}}

\subsection{Broader Context}

These House actions sit within a wider policy landscape. The 2015 SPACE Act and the 2017 American Space Commerce Free Enterprise Act were particularly notable for their implications under international law, as they effectively challenged the long-standing framework of the 1967 Outer Space Treaty, which established that outer space is “not subject to national appropriation by claim of sovereignty”.\footnote{United Nations Office for Outer Space Affairs. (n.d.). \emph{The Outer Space Treaty}. \url{https://www.unoosa.org/oosa/en/ourwork/spacelaw/treaties/introouterspacetreaty.html}}

Most recently, President Trump’s August 2025 Executive Order on “Enabling Competition in the Commercial Space Industry” directed the streamlining of commercial licences and permits and called on the Secretary of Commerce to propose a new process for mission authorisations for novel private-sector space activities.\footnote{The White House. (2025, August 13). \emph{Enabling competition in the commercial space industry} [Presidential action]. \url{https://www.whitehouse.gov/presidential-actions/2025/08/enabling-competition-in-the-commercial-space-industry/}}

The Belfer Centre has noted that this executive order acknowledges a growing class of commercially provided in-space services, including in-space servicing and assembly, orbital debris removal, in-space resource utilisation, and in-space manufacturing, that lack clear governance under existing U.S. regulation.\footnote{Belfer Center for Science and International Affairs. (2025). \emph{Explainer: Novel space activities in the executive order on enabling competition in the commercial space industry}. Harvard Kennedy School. \url{https://www.belfercenter.org/research-analysis/explainer-novel-space-activities-executive-order-enabling-competition-commercial}}
